%%
% Please see https://bitbucket.org/rivanvx/beamer/wiki/Home for obtaining beamer.
%%
%\documentclass[11pt,aspectratio=169]{beamer}
%\documentclass[11pt,aspectratio=169,xcolor={dvipsnames}]{beamer}

\documentclass[11pt,aspectratio=169,xcolor={dvipsnames},hyperref={pdftex,pdfpagemode=UseNone,hidelinks,pdfdisplaydoctitle=true},usepdftitle=false]{beamer}
\usepackage{presentation}


\title{Lecture 16: Q-Theory of Investment}
\author{Juan Herre\~{n}o\\UCSD}
\date{2026}

% Enter presentation title to populate PDF metadata:
\newcommand{\Cov}{\text{Cov}_\lambda}
\newcommand{\El}{\mathbb{E}_\lambda}
\newcommand{\mm}{\text{{\fontfamily{qzc}\selectfont m}}}
\usepackage{cancel}
\usepackage{booktabs}
\begin{document}

\maketitle


\begin{frame}{The Neoclassical Model}
Time is continuous (same framework you saw in Econ 210B). Firms:
\pause
\begin{itemize}
\item Discounts future profits at rate $r$
\pause
\item Produces a homogeneous good
\pause
\item Perfect competition in input and output markets
\pause
\item Production Function is CRS
\pause
\item Hires labor, buys capital
\end{itemize}
Problem of a firm that takes input prices as given
\end{frame}

\begin{frame}{Profits}
Profits are given by revenues minus the wage bill minus capital expenditures
\pause
\begin{align}
\pi_{t} = A_t F(k_t,n_t) - w_t n_t - i_t
\end{align}
\pause
Capital adjusts given the law of motion
\pause
\begin{align}
\dot{k}_t = i_t
\end{align}
\pause
\begin{itemize}
\item We are assuming zero depreciation just for simplicity
\pause
\item Abstracting from taxes, bonus depreciation, etc.
\pause
\item Perfect foresight
\end{itemize}
\end{frame}

\begin{frame}{The firm's problem}
Objective
\pause
\begin{align}
\max_{i_t, n_t} V = \int_0^{\infty} e^{-rt} \pi_t dt
\end{align}
\pause
Subject to
\pause
\begin{align}
\pi_{t} &= A_t F(k_t,n_t) - w_t n_t - i_t\\
\dot{k}_t& = i_t
\end{align}
\pause
You are experts in solving these problems
\pause
\begin{align}
\mathcal{H} = e^{-rt} \left(A F(k_t,n_t) - w_t N_t -  i_t \right) + \lambda_t i_t.
\end{align}
\pause
Question: What is the meaning of $\lambda_t$?
\end{frame}

\begin{frame}{$\lambda$ vs. $q$}
\pause
$\lambda_t$ is the marginal value of one unit of investment at time $t$ to time-zero firm value
\begin{itemize}
\pause
\item It is the \textit{present value} marginal value of investment
\end{itemize}
\pause
We will introduce a new concept, the current marginal value $q_t$
\begin{align}
\pause
q_t = \lambda_t e^{rt}
\end{align}
\end{frame}

\begin{frame}{Optimality Conditions $\lambda$}
\begin{align}
\mathcal{H} = e^{-rt} \left(A F(k_t,n_t) - w_t n_t -  I_t \right) + \lambda_t i_t.
\end{align}
\pause
Optimality conditions:
\pause
\begin{align}
A F_n &= w_t\\
e^{-rt} &= \lambda_t \\
e^{-rt} A_t F_k &= -\dot{\lambda}_t \\
\lim_{t\rightarrow \infty} \lambda_t k_t &= 0.
\end{align}
\end{frame}


\begin{frame}{Optimality Conditions $q$}
\begin{align}
\mathcal{H} = e^{-rt} \left(A F(k_t,n_t) - w_t n_t -  I_t \right) + \lambda_t i_t.
\end{align}
\pause
Optimality conditions:
\begin{align}
A F_n &= w_t\\
q_t &= 1 \\
A F_k &= - \dot{q}_t + r q_t\\
\lim_{t\rightarrow \infty} e^{-rt} q_t k_t &= 0.
\end{align}
\end{frame}

\begin{frame}{Zoom in the key result}
$$q_t = 1$$
\pause
What is the meaning of this equation?
\pause
\begin{itemize}
\item The marginal value of investing is equal to $q_t$
\pause
\item The marginal cost of investing is losing 1 unit of profits
\pause
\item Optimality requires to invest until marginal benefit equals marginal cost
\end{itemize}
\end{frame}


\begin{frame}{Implications of the key result}
If $q_t = 1$, then $\dot{q}_t = 0$\\
\pause
So equation $A F_k = - \dot{q}_t + r q_t$\\
\pause
Transforms into
\pause
$$A F_k = r $$
It is as if firms are solving a static problem
\pause
\begin{itemize}
\item Equate marginal product of capital to $r$ at each point in time
\pause
\item There is no role for future variables
\pause
\item There is no role for past variables
\end{itemize}
\end{frame}

\begin{frame}{The economics of the basic problem}
Imagine a firm starts with $k_0 < k^*$ , for $k^*$ consistent with $AF_k^* = r$
\pause
\begin{itemize}
\item Capital will jump immediately to $k^*$
\pause
\item Investment is not well-defined at the time of adjustment
\pause
\item It will be zero otherwise
\end{itemize}
\pause
Not having a properly defined investment function is obviously problematic. Economics:
\begin{itemize}
\pause
\item Capital should \textit{jump} instantaneously
\pause
\item Therefore, investment needs to be infinite for an instant
\end{itemize}
\end{frame}


\begin{frame}{The economics of the basic problem}
Imagine $A$ increases to $A_H$ and goes down to $A$ in the following instant
\begin{itemize}
\item Firms will infinitely invest in the present
\pause
\item Firms will infinitely disinvest in the following instant
\end{itemize}
Against any possible intuition. $r = A F_k$
\end{frame}

\begin{frame}{Away from continuous time for a second}
You may think this are problems that arise only in the continuous time limit
\end{frame}


\begin{frame}{Away from continuous time for a second}
\begin{itemize}
\item To further illustrate this point imagine we are in discrete time, each period is a year.
\pause
\item The same equation applies
$$A F_k = r $$
\pause
\item Assume that $F(K,L) = K^{\alpha}L^{1-\alpha}$. Therefore:
$$1+r_t = 1+ A_t \alpha K_t^{\alpha-1}L_t^{1-\alpha}$$
\vspace{-0.5cm}
\pause
\item Make a log-linear approximation. Hatted variables are log changes:
$$\hat{r}_t = \frac{r}{1+r}\left(\hat{a}_t - (1-\alpha) \hat{k}_t + (1-\alpha) \hat{l}_t\right)$$
\vspace{-0.5cm}
\pause
\item where $\hat{r_t} = \log \frac{1+r_t}{1+r}$, solve for $\hat{k}$
$$\hat{k}_t = - \hat{r}_t \left(\frac{1+r}{(1-\alpha)r}\right) + \frac{\hat{a}_t}{1-\alpha} + \hat{l}_t$$
\vspace{-0.5cm}
\end{itemize}
\end{frame}

\begin{frame}{Away from continuous time for a second}
$$\hat{k}_t = - \hat{r}_t \left(\frac{1+r}{(1-\alpha)r}\right) + \frac{\hat{a}_t}{1-\alpha} + \hat{l}_t$$
\begin{itemize}
\pause
\item Assume an exogenous decrease of 1\% in interest rates.
\pause
\item Capital would have to increase 31.5\%  
\pause
\item Including reasonable depreciation would change this number to 14\%.
\pause
\item Letting labor increase would further increase this number
\pause
\item Assume 100\% of GDP could be transformed to capital.
\pause
\item Capital-output ratios are between 2 and 4 (depending on land and housing)
\pause
\item To increase capital by 31\%, it would take 61\%-124\% of GDP
\item With $\delta$: to increase capital by 14\%, it would take 28-56\% of GDP
\end{itemize}
\end{frame}

\begin{frame}{The economics of the basic problem}
Imagine firms notice that $A$ will increase at time $\tau$ ($A_{\tau} > A$)
\begin{itemize}
\item Firms do nothing up until period $\tau$
\item invest $+ \infty$ at time $\tau$
\end{itemize}
No role for news about the future
\end{frame}

\begin{frame}{Idea: convex adjustment costs}
Investing takes time and resources
\pause
\begin{itemize}
\item Installing capital is costly
\item Planning the right level of capital is costly
\item scrapping old capital is costly
\end{itemize}
\pause
Convex costs
\begin{align}
\phi(i,k) = \frac{\varphi}{2} \left(\frac{i}{k}\right)^2 k
\end{align}
Lucas (1967) is an early reference for convex costs on investment
\begin{itemize}
\item Note that defined in this way $\phi(i,k)$ is homogeneous of degree 1
\end{itemize}
\end{frame}

\begin{frame}{New firm's problem}
Objective
\begin{align}
\max_{i_t, n_t} V = \int_0^{\infty} e^{-rt} \pi_t dt
\end{align}
Subject to
\begin{align}
\pi_{t} &= A_t F(k_t,n_t) - w_t n_t - i_t - \frac{\varphi}{2} \left(\frac{i}{k}\right)^2 k \\
\dot{k}_t& = i_t
\end{align}
You are experts in solving these problems
\pause
\begin{align}
\mathcal{H} = e^{-rt} \left(A F(k_t,n_t) - w_t N_t -  i_t - \frac{\varphi}{2} \left(\frac{i}{k}\right)^2 k \right) + \lambda_t i_t.
\end{align}
\pause
Question: What is the meaning of $\lambda_t$?
\end{frame}


\begin{frame}{Optimality Conditions}
\begin{align}
A F_N &= w_t\\
1 + \varphi \frac{i_t}{k_t}  &= q_t\\
-\dot{q}_t + r q_t &= AF_k + \frac{\varphi}{2}\left(\frac{i}{k}\right)^2
\end{align}
\end{frame}

\begin{frame}{Zoom in the important economics}
$$1 + \varphi \frac{i_t}{k_t}  = q_t$$
\begin{itemize}
\item Rewrite the equation as
$$\frac{i_t}{k_t} = h(q_t)$$
\item The investment rate is \textbf{only} a function of $q_t$
\pause 
\item $q$ is a \textit{sufficient statistic}
\end{itemize}
\end{frame}

\begin{frame}{Two benefits for having more capital}
\begin{itemize}
\item You can solve the differential equation
$$-\dot{q}_t + r q_t = AF_k +\frac{\varphi}{2} \left(\frac{i}{k}\right)^2$$
\item And find that
$$q_0 = \int_0^{\infty} e^{-rt} \left[ A_t F_{kt} +  \frac{\varphi}{2}  \left(\frac{i_t}{k_t}\right)^2 \right]dt$$
\item Investment brings a sequence of marginal products
\pause
\item Reductions in future costs of adjustment
\pause
\item that is the marginal value of investment.
\end{itemize}
\end{frame}

\begin{frame}{The mechanics of the q-Theory}
\begin{itemize}
\item Imagine there are news about future profitability
\pause
\item This increases $A$ for instance
\item Increasing $q$ by
$$q_0 = \int_0^{\infty} e^{-rt} \left[ A_tF_{kt} +  \frac{\varphi}{2} \left(\frac{i_t}{k_t}\right)^2 \right]dt$$
\pause
\item a higher $q$ calls for higher investment
$$\frac{i_t}{k_t}  = \frac{q_t -1}{\varphi} $$
\item Higher $k$ brings down future MPK
\item Lowering $q$
\end{itemize}
\end{frame}

\begin{frame}{When to invest}
$$\frac{i_t}{k_t}  = \frac{q_t -1}{\varphi} $$
\begin{itemize}
\pause
\item Investment will be positive as long as $q>1$
\pause
\item How much to invest depends on $\varphi$
\pause
\item What happens in the limit $\varphi \rightarrow 0$?
\end{itemize}
\end{frame}

\begin{frame}{How to implement the Q Theory}
\begin{itemize}
\item So far we have a theory that ties investment to $q$
\pause
\item We do not have good measures for marginal $q$ in the data
\pause
\item It is a similar measurement problem to that of having measures of marginal cost!
\pause
\item We do have measures of \textit{average} $q$. Also called \textit{Tobin's} Q
\item To which we will refer to as $Q$.
\end{itemize}
\end{frame}

\begin{frame}{Tobin's Q}
\begin{itemize}
\item Tobin (1969) argued that firms should invest if 
$$Q = \frac{\text{Value of the Firm Capital}}{\text{Replacement Cost of the Capital Stock}}$$
\item Capital is valued more within the firm, than outside of the firm.
\pause
\item Our theory says that investment should be higher if $q>1$
\pause
\item Under what conditions Tobin's claim is sustained in our standard model?
\end{itemize}
\end{frame}


\begin{frame}{How to implement the Q Theory}
\begin{itemize}
\item $q$ is the marginal value of one unit of capital
\pause
\item Average $Q$ is the average value of a unit of capital
$$Q = \frac{\text{Value of the Firm Capital}}{\text{Replacement Cost of the Capital Stock}}$$
\pause
\item Can we learn about $q$ with information about $Q$?
\end{itemize}
\end{frame}


\begin{frame}{Hayashi 1982}
\begin{itemize}
\item Under a number of conditions, there is an equivalent result.
\begin{itemize}
\item Firms have CRS technology
\item The adjustment cost technology is CRS
\item Perfect competition
\item Efficient Financial Markets
\end{itemize}
\item Then $q_t = Q_t$
\item Hayashi (1982) presented this result in continuous time.
\item I will show you the result in discrete time, which for me is more intuitive
\end{itemize}
\end{frame}


\begin{frame}{Marginal $q$ and Average $Q$ in discrete time}
\begin{itemize}
\item Time is discrete and the problem is very similar to what we covered before
\pause
\item Profits are given by
$$\pi_t = AF(K_t,L_t) - W_t L_t - I_t - \psi \left(\frac{I_t}{K_t}\right) K_t$$
\vspace{-0.5cm}
\item Firms maximize their value
$$V_0 = \sum_{t=0}^{\infty} \left(\frac{1}{1+r}\right)^{t} \pi_t$$
\vspace{-0.5cm}
\pause
\item Subject to the law of motion of capital
$$K_{t+1} = K_t(1-\delta) + I_t$$
\vspace{-0.5cm}
\item $q_t$ is the Lagrange multiplier
\end{itemize}
\end{frame}

\begin{frame}{Marginal $q$ and Average $Q$ in discrete time}
\begin{itemize}
\item The problem has very similar first order conditions
$$AF_L(K_t,L_t) = W_t$$
$$q_t = 1 + \psi' \left(\frac{I_t}{K_t}\right)$$
$$q_t = \frac{1-\delta}{1+r} q_{t+1} + \frac{1}{1+r} \left(A_{t+1} F_k(K_{t+1},L_{t+1}) - \psi \left(\frac{I_{t+1}}{K_{t+1}}\right) - \psi' \left(\frac{I_{t+1}}{K_{t+1}}\right) \frac{I_{t+1}}{K_{t+1}}\right)$$
\end{itemize}
\end{frame}

\begin{frame}{Marginal $q$ and Average $Q$ in discrete time}
$$q_t = \frac{1-\delta}{1+r} q_{t+1} + \frac{1}{1+r} \left(A_{t+1} F_k(K_{t+1},L_{t+1}) - \psi \left(\frac{I_{t+1}}{K_{t+1}}\right) - \psi' \left(\frac{I_{t+1}}{K_{t+1}}\right) \frac{I_{t+1}}{K_{t+1}}\right)$$
\vspace{-0.5cm}
\begin{itemize}
\pause
\item And multiply by $K_{t+1}$
$$q_t K_{t+1} = \frac{1-\delta}{1+r} q_{t+1}K_{t+1} + \frac{1}{1+r} \left(A_{t+1} F_k(K_{t+1},L_{t+1})K_{t+1} - \psi \left(\frac{I_{t+1}}{K_{t+1}}\right)K_{t+1}  - \psi' \left(\frac{I_{t+1}}{K_{t+1}}\right) I_{t+1} \right)$$
\vspace{-0.5cm}
\pause
\item Now use $q_{t+1} = 1 + \psi' \left(\frac{I_{t+1}}{K_{t+1}}\right)$, to replace $\psi'$
$$q_t K_{t+1} = \frac{1-\delta}{1+r} q_{t+1}K_{t+1} + \frac{1}{1+r} \left(A_{t+1} F_k(K_{t+1},L_{t+1})K_{t+1} - \psi \left(\frac{I_{t+1}}{K_{t+1}}\right)K_{t+1}  - (q_{t+1} - 1) I_{t+1} \right)$$
\vspace{-0.5cm}
\pause
\item Use the law of motion of capital in period $t+1$ to replace $I_{t+1}$
$$q_t K_{t+1} = \frac{1-\delta}{1+r} q_{t+1}K_{t+1} + \frac{1}{1+r} \left(A_{t+1} F_k(K_{t+1},L_{t+1})K_{t+1} - \psi \left(\frac{I_{t+1}}{K_{t+1}}\right)K_{t+1}  - (q_{t+1} - 1)(K_{t+2} - K_{t+1}(1-\delta))\right)$$
\vspace{-0.5cm}
\end{itemize}
\end{frame}

\begin{frame}{Marginal $q$ and Average $Q$ in discrete time}
\begin{itemize}
\item Rearrange
$$q_t K_{t+1} =   \frac{1}{1+r}q_{t+1}K_{t+2} + \frac{1}{1+r} \left(A_{t+1} F_k(K_{t+1},L_{t+1})K_{t+1} - \psi \left(\frac{I_{t+1}}{K_{t+1}}\right)K_{t+1}  - I_{t+1} \right)$$
\vspace{-0.5cm}
\item Now use Euler's Theorem $F(K,L) = F_k(K,L) K + F_L(K,L) L$
\item And, perfect competition in the labor market $W = A F_L (K,L)$
$$q_t K_{t+1} =   \frac{1}{1+r}q_{t+1}K_{t+2} + \frac{1}{1+r} \left(A_{t+1} F(K_{t+1},L_{t+1}) - W_{t+1} L_{t+1} - \psi \left(\frac{I_{t+1}}{K_{t+1}}\right)K_{t+1}  - I_{t+1} \right)$$
\vspace{-0.5cm}
\item Notice that the term in parenthesis is just profits tomorrow
$$q_t K_{t+1} =   \frac{1}{1+r}q_{t+1}K_{t+2} + \frac{1}{1+r} \pi_{t+1}$$
\end{itemize}
\end{frame}

\begin{frame}{Marginal $q$ and Average $Q$ in discrete time}
\begin{itemize}
\item Iterate forward and use the TVC to get rid of the terminal term
$$q_t K_{t+1} =  \sum_{j=1}^{\infty} \left(\frac{1}{1+r}\right)^t \pi_{t+1}$$
\vspace{-0.5cm}
\item Use efficiency of financial markets:
$$q_t K_{t+1} = V^e_{t}$$
\pause
\item $V_t^e$ is the ``ex-dividend'' (after paying today's dividend) value of the firm equity
\item Divide by the stock of capital at the beginning of tomorrow's period
$$q_t =\frac{ V^e_{t}}{K_{t+1}}$$
\vspace{-0.5cm}
\pause
\item Realize that the RHS is the average $Q$
$$q_t = Q_t$$
\end{itemize}
\end{frame}

\begin{frame}{Relevance of this Result}
\begin{itemize}
\item $q_t = Q_t$. So?
\pause
\item This result makes the Q-theory operational
\pause
\item $Q$ is measurable
\pause
\item We can now measure $Q$ at different levels of aggregation
\begin{itemize}
\item Firm level
\item Industry level
\item National level
\end{itemize}
\item  And test the theory
\pause
\item Why do we want to test this theory?
\begin{itemize}
\item It is our benchmark model
\item Many educated guesses come from this framework
\item The investment effects of tax reforms
\item The investment effects of monetary policy
\item And so on...
\end{itemize}
\end{itemize}
\end{frame}



\end{document}


